% =====================================================================
%  preamble.tex
%  Configuración global del documento: paquetes, idioma, estilo,
%  cajas destacadas y bibliografía.
%
%  Motor recomendado: pdflatex (compila también con xelatex/lualatex).
%  Bibliografía: BibTeX + natbib (estilo autor-año).
% =====================================================================

% ---------------------------------------------------------------------
%  Codificación y tipografía
% ---------------------------------------------------------------------
\usepackage[T1]{fontenc}      % Codificación de fuente moderna (acentos)
\usepackage[utf8]{inputenc}   % El fichero fuente está en UTF-8
% Latin Modern (fuentes vectoriales con buena cobertura de acentos).
% Se carga sólo si está disponible; en distribuciones completas y en
% Overleaf lo está. Si falta, se usa la Computer Modern por defecto.
\IfFileExists{lmodern.sty}{\usepackage{lmodern}}{}
% microtype: protrusion siempre; expansion=false evita el error fatal
% «auto expansion is only possible with scalable fonts» cuando, por
% falta de lmodern, se recurre a fuentes de mapa de bits.
\usepackage[protrusion=true, expansion=false]{microtype}
\usepackage{textcomp}         % Símbolos adicionales

% ---------------------------------------------------------------------
%  Idioma español
%  Se usa babel moderno con \babelprovide para no depender del antiguo
%  spanish.ldf. Aporta «Capítulo», «Índice», fechas y guiones en español.
% ---------------------------------------------------------------------
\usepackage{babel}
\babelprovide[import, main]{spanish}

% ---------------------------------------------------------------------
%  Geometría de página
% ---------------------------------------------------------------------
\usepackage[
  a4paper,
  top=2.6cm,
  bottom=2.6cm,
  inner=2.8cm,
  outer=2.4cm,
  headsep=14pt,
]{geometry}

% ---------------------------------------------------------------------
%  Color: paleta del manual
% ---------------------------------------------------------------------
\usepackage{xcolor}
\definecolor{primario}{HTML}{1F4E5F}   % Azul petróleo (títulos, filetes)
\definecolor{secundario}{HTML}{8E3B46} % Granate (acentos)
\definecolor{aviso}{HTML}{B23A48}      % Rojo aviso
\definecolor{consejo}{HTML}{2E7D5B}    % Verde consejo
\definecolor{regla}{HTML}{B07D2B}      % Ámbar regla de oro
\definecolor{grisclaro}{HTML}{F4F4F1}  % Fondo suave de cajas

% ---------------------------------------------------------------------
%  Listas, tablas y figuras
% ---------------------------------------------------------------------
\usepackage{enumitem}         % Control fino de listas
\setlist{itemsep=2pt, topsep=4pt, parsep=0pt}
\usepackage{booktabs}         % Filetes de tabla de calidad
\usepackage{tabularx}         % Tablas de ancho fijo
\usepackage{array}            % Mejoras en columnas
\usepackage{ragged2e}         % \RaggedRight en celdas
\usepackage{graphicx}         % Imágenes (carpeta figures/)
\graphicspath{{figures/}}

% ---------------------------------------------------------------------
%  Comillas (estilo español: «...»)
% ---------------------------------------------------------------------
\usepackage{csquotes}

% ---------------------------------------------------------------------
%  Cajas destacadas (avisos, consejos, definiciones)
% ---------------------------------------------------------------------
\usepackage[most]{tcolorbox}

% Caja genérica con título e icono textual
\newtcolorbox{cajadestacada}[2][]{%
  enhanced, breakable,
  colback=grisclaro, colframe=#2,
  boxrule=0.4pt, left=10pt, right=10pt, top=8pt, bottom=8pt,
  borderline west={3pt}{0pt}{#2},
  arc=2pt, fonttitle=\bfseries\color{#2},
  #1
}

% Aviso importante (rojo)
\newtcolorbox{aviso}{%
  enhanced, breakable,
  colback=aviso!6, colframe=aviso,
  boxrule=0.4pt, left=10pt, right=10pt, top=8pt, bottom=8pt,
  borderline west={3pt}{0pt}{aviso}, arc=2pt,
}

% Consejo (verde)
\newtcolorbox{consejo}{%
  enhanced, breakable,
  colback=consejo!6, colframe=consejo,
  boxrule=0.4pt, left=10pt, right=10pt, top=8pt, bottom=8pt,
  borderline west={3pt}{0pt}{consejo}, arc=2pt,
}

% Regla de oro (ámbar)
\newtcolorbox{reglaoro}{%
  enhanced, breakable,
  colback=regla!8, colframe=regla,
  boxrule=0.4pt, left=10pt, right=10pt, top=8pt, bottom=8pt,
  borderline west={3pt}{0pt}{regla}, arc=2pt,
}

% Definición / concepto (azul primario)
\newtcolorbox{definicion}{%
  enhanced, breakable,
  colback=primario!5, colframe=primario,
  boxrule=0.4pt, left=10pt, right=10pt, top=8pt, bottom=8pt,
  borderline west={3pt}{0pt}{primario}, arc=2pt,
}

% ---------------------------------------------------------------------
%  Estilo de títulos
% ---------------------------------------------------------------------
\usepackage{titlesec}

\titleformat{\chapter}[display]
  {\normalfont\sffamily\bfseries\color{primario}}
  {\filright\large\sffamily Capítulo \thechapter}
  {6pt}
  {\Huge}
  [\vspace{2pt}{\color{primario}\titlerule[1pt]}]
\titlespacing*{\chapter}{0pt}{10pt}{28pt}

\titleformat{\section}
  {\normalfont\sffamily\large\bfseries\color{primario}}
  {\thesection}{0.8em}{}
\titleformat{\subsection}
  {\normalfont\sffamily\bfseries\color{secundario}}
  {\thesubsection}{0.6em}{}

% ---------------------------------------------------------------------
%  Encabezados y pies de página
% ---------------------------------------------------------------------
\usepackage{fancyhdr}
\pagestyle{fancy}
\fancyhf{}
\renewcommand{\headrulewidth}{0.4pt}
\renewcommand{\footrulewidth}{0pt}
\fancyhead[LE]{\small\nouppercase{\leftmark}}
\fancyhead[RO]{\small\nouppercase{\rightmark}}
\fancyfoot[C]{\thepage}
% Estilo «plain» (primera página de capítulo) sólo con el número
\fancypagestyle{plain}{%
  \fancyhf{}\fancyfoot[C]{\thepage}%
  \renewcommand{\headrulewidth}{0pt}}

% ---------------------------------------------------------------------
%  Espaciado e interlineado
% ---------------------------------------------------------------------
\usepackage{setspace}
\setstretch{1.25}
\setlength{\parskip}{1.0em}
\setlength{\parindent}{0pt}
\usepackage{emptypage}        % Páginas en blanco realmente en blanco

% ---------------------------------------------------------------------
%  Bibliografía (autor-año con natbib + BibTeX)
% ---------------------------------------------------------------------
\usepackage[round, authoryear]{natbib}
\bibliographystyle{plainnat}
\setcitestyle{aysep={,}}

% ---------------------------------------------------------------------
%  Hiperenlaces y marcadores (cargar casi al final)
% ---------------------------------------------------------------------
\usepackage[
  colorlinks=true,
  linkcolor=primario,
  citecolor=secundario,
  urlcolor=secundario,
  pdftitle={Defensa Personal para Mujeres. Manual del Estudiante},
  pdfauthor={Artes Marciales Bilbao},
  pdfsubject={Manual de defensa personal y autoprotección},
]{hyperref}
\usepackage{bookmark}
\urlstyle{same}

% ---------------------------------------------------------------------
%  Comandos auxiliares
% ---------------------------------------------------------------------
% Teléfono/recurso destacado en línea
\newcommand{\recurso}[1]{\textbf{\color{primario}#1}}
